\scanpage{42}
\mspage{16}
% archive foliation: 44
ne donne quelque chose que si on la complète par la prescription que l'on prend une métrique ``distinguée'' sur $X^{or}$ qui soit stable par $\sigma$.

Considérons maintenant les seuls cas nouveaux --- ceux où $X$ est 0-connexe et $\partial X\ne\varnothing$, correspondant à un $X^{or}=\mathcal X$ connexe, munie d'une anti-involution $\sigma$ telle que $\mathcal X^\sigma\ne\varnothing$. Les seuls cas problématiques pour l'existence d'une métrique canonique sur $X$, sont les suivants :

a') $\mathcal X\simeq\mathbb C$. \deleted{Les \unsure{anti.}} \inserted{OPS} $\mathcal X=\mathbb C$. Les anti-aut. de $\mathbb C$ sont
\[
z\longmapsto a\overline z+b\qquad(a\in\mathbb C^*,\ b\in\mathbb C).
\]
Si on prend une anti-involution à pts fixes, OPS que 0 est pt fixe, donc $b=0$; donc \deleted{l'application est $z\longmapsto a\overline z$}
\[
\sigma z=a\overline z\qquad\text{avec }a\overline a=1\quad\text{i.e }|a|=1\quad\text{i.e }a\in\mathbb U
\]
(pour avoir $\sigma^2=\mathrm{id}$) soit
\deleted{l'ens des pts fixes est}
\[
\sigma z=\alpha^2\overline z=\frac\alpha{\overline\alpha}\overline z\qquad\text{si }a=\alpha^2,\quad\alpha\in\mathbb U,
\]
et l'ens des pts fixes est
\[
\mathbb C^\sigma=\alpha.\mathbb R.
\]
\margin{symétrie p.r. à la droite $\alpha.\mathbb R$}
